% 无线电仪器与仪表
% 无线电仪器与仪表.tex

\documentclass[12pt,UTF8]{ctexbook}

% 设置纸张信息。
\input{../../../common/page}

% 设置字体，并解决显示难检字问题。
\xeCJKsetup{AutoFallBack=true}
% 注意：ctexbook类已默认设置SimSun为CJKrmdefault
% 以下设置用于确保字体回退和扩展字体可用
% 仅设置扩展字体以避免与默认设置冲突
\setCJKfamilyfont{hei}{SimHei}
\setCJKfamilyfont{kai}{KaiTi}

% 目录 chapter 级别加点（.）。
\usepackage{titletoc}
\titlecontents{chapter}[0pt]{\vspace{3mm}\bf\addvspace{2pt}\filright}{\contentspush{\thecontentslabel\hspace{0.8em}}}{}{\titlerule*[8pt]{.}\contentspage}

% 支持目录点击跳转
\usepackage[colorlinks,linkcolor=blue,citecolor=blue,urlcolor=blue]{hyperref}

% 设置 part 和 chapter 标题格式。
\ctexset{
	part/name= {第,卷},
	part/number={\chinese{part}},
	chapter/name={第,篇},
	chapter/number={\arabic{chapter}}
}

% 图片相关设置。
\usepackage{graphicx}
\graphicspath{{Images/}}

% 设置署名格式。
\newenvironment{shuming}{\hfill\zihao{4}}

% 注脚每页重新编号，避免编号过大。
\usepackage[perpage]{footmisc}

% 设置古文原文格式。
\newenvironment{yuanwen}{\bfseries\zihao{4}}

% 设置段落间距
\setlength{\parskip}{0.8em plus 0.2em minus 0.1em}

% 列表格式支持，列表项向右偏移。
\usepackage{enumitem}

\begin{document}

\frontmatter

\title{无线电仪器与仪表}
\author{}
\date{}
\maketitle

\tableofcontents

\mainmatter

\part{无线电测量基础}

\chapter{无线电测量概述}
\section{无线电测量的基本概念}
\section{无线电测量的重要性}
\section{无线电测量的发展历程}

\chapter{测量误差与精度}
\section{测量误差的分类}
\section{测量精度的评定}
\section{测量不确定度分析}

\chapter{测量单位与标准}
\section{无线电测量单位体系}
\section{国际测量标准}
\section{国家标准与校准}

\part{基本无线电测量仪器}

\chapter{电压测量仪器}
\section{模拟电压表}
\section{数字电压表}
\section{高频电压表}
\section{峰值电压表}

\chapter{电流测量仪器}
\section{模拟电流表}
\section{数字电流表}
\section{高频电流探头}
\section{电流钳形表}

\chapter{功率测量仪器}
\section{功率计的基本原理}
\section{射频功率计}
\section{微波功率计}
\section{功率传感器}

\chapter{频率测量仪器}
\section{频率计的基本原理}
\section{数字频率计}
\section{频率计数器}
\section{微波频率计}

\part{信号分析仪器}

\chapter{示波器}
\section{模拟示波器}
\section{数字示波器}
\section{存储示波器}
\section{混合信号示波器}

\chapter{频谱分析仪}
\section{频谱分析仪的基本原理}
\section{扫频式频谱分析仪}
\section{实时频谱分析仪}
\section{矢量信号分析仪}

\chapter{网络分析仪}
\section{网络分析仪的基本原理}
\section{标量网络分析仪}
\section{矢量网络分析仪}
\section{微波网络分析仪}

\chapter{信号发生器}
\section{函数信号发生器}
\section{射频信号发生器}
\section{微波信号发生器}
\section{任意波形发生器}

\part{专用无线电测量仪器}

\chapter{天线测量仪器}
\section{天线分析仪}
\section{天线测试系统}
\section{近场测试系统}
\section{远场测试系统}

\chapter{电磁兼容测量仪器}
\section{EMI接收机}
\section{电磁场探头}
\section{传导发射测试系统}
\section{辐射发射测试系统}

\chapter{通信测试仪器}
\section{误码率测试仪}
\section{协议分析仪}
\section{无线通信测试仪}
\section{移动通信测试仪}

\section{学生信号源（J2465型）}

J2465型学生信号源是一种专门为教学实验设计的无线电仪器，具有以下特点：

\subsection{主要功能}
\begin{itemize}[leftmargin=2cm]
	\item \textbf{信号产生}：提供多种波形信号输出
	\item \textbf{教学用途}：专门为无线电教学实验设计
	\item \textbf{操作简便}：适合学生实验操作
	\item \textbf{安全可靠}：符合教学仪器安全标准
\end{itemize}

\subsection{技术参数}
J2465型学生信号源是具有高频信号和音频信号输出的简易信号源，主要技术指标如下：

\begin{itemize}[leftmargin=2cm]
	\item \textbf{高频信号频率范围}：400kHz～1700kHz，分两个频段
	\begin{itemize}[leftmargin=1cm]
		\item 频段Ⅰ：1500kHz～1700kHz
		\item 频段Ⅱ：400kHz～580kHz
	\end{itemize}
	\item \textbf{低频信号频率范围}：音频信号输出
	\item \textbf{输出波形}：正弦波为主，适合教学实验
	\item \textbf{输出幅度}：可调输出电平，适合学生实验
	\item \textbf{调制功能}：支持调幅（AM）调制
	\item \textbf{工作电源}：直流电源供电，符合教学安全要求
	\item \textbf{输出接口}：标准BNC接口或接线柱输出
\end{itemize}

\subsection{应用领域}
\begin{itemize}[leftmargin=2cm]
	\item \textbf{无线电基础教学实验}
	\item \textbf{通信原理课程实验}
	\item \textbf{电子技术实训}
	\item \textbf{学生创新项目开发}
\end{itemize}

\subsection{使用方法}
J2465学生信号源的使用方法相对简单，适合教学实验：

\begin{itemize}[leftmargin=2cm]
	\item \textbf{频率选择}：通过波段开关选择高频或低频输出
	\item \textbf{频率调节}：使用频率旋钮在选定频段内调节具体频率
	\item \textbf{幅度调节}：调节输出信号的电平大小
	\item \textbf{调制设置}：根据需要设置调幅调制深度
	\item \textbf{输出连接}：通过标准接口连接到被测设备
\end{itemize}

\subsection{注意事项}
使用J2465学生信号源时需要注意以下事项：

\begin{itemize}[leftmargin=2cm]
	\item \textbf{电源安全}：确保使用正确的直流电源电压
	\item \textbf{接地处理}：注意信号源接地，避免干扰问题
	\item \textbf{频率稳定}：预热一段时间后频率更加稳定
	\item \textbf{输出保护}：避免输出短路或过载
	\item \textbf{教学指导}：建议在教师指导下使用
\end{itemize}

\subsection{典型应用实验}
J2465学生信号源适用于以下典型教学实验：

\begin{itemize}[leftmargin=2cm]
	\item \textbf{收音机调试}：用于中波收音机的调试和维修
	\item \textbf{电路测试}：测试放大电路、滤波电路等
	\item \textbf{信号观察}：配合示波器观察信号波形
	\item \textbf{无线实验}：进行简单的无线发射实验
\begin{itemize}[leftmargin=1cm]
	\item 调幅无线发射实验
	\item 无线信号传播特性实验
	\item 收音机接收实验
	\item 天线效应实验
\end{itemize}
\end{itemize}

J2465作为教学用信号源，在无线电教育领域具有重要地位，是培养学生无线电技术基础能力的重要工具。该信号源虽然功能相对简单，但完全满足基础教学实验的需求，具有操作简便、安全可靠的特点。

\subsection{无线实验详细操作方法}
使用J2465学生信号源进行无线实验的具体步骤如下：

\subsubsection{调幅无线发射实验}
\begin{enumerate}[leftmargin=2cm]
	\item \textbf{设备准备}：J2465信号源、收音机、连接线、天线
	\item \textbf{频率设置}：将信号源频率设置在收音机中波频段（如1000kHz）
	\item \textbf{调制设置}：开启调幅功能，设置适当的调制深度
	\item \textbf{天线连接}：将天线连接到信号源的高频输出端
	\item \textbf{收音机调谐}：将收音机调谐到对应的频率
	\item \textbf{信号接收}：在收音机中应能听到调制信号的声音
\end{enumerate}

\subsubsection{无线信号传播特性实验}
\begin{enumerate}[leftmargin=2cm]
	\item \textbf{距离测试}：在不同距离测试信号强度变化
	\item \textbf{方向性测试}：测试信号在不同方向的传播特性
	\item \textbf{障碍物影响}：测试墙壁、金属等障碍物对信号的影响
	\item \textbf{多径效应}：观察信号反射造成的多径效应
\end{enumerate}

\subsubsection{收音机接收实验}
\begin{enumerate}[leftmargin=2cm]
	\item \textbf{灵敏度测试}：测试收音机的最小可接收信号强度
	\item \textbf{选择性测试}：测试收音机对邻近频率的抑制能力
	\item \textbf{信噪比测试}：测量接收信号的信噪比
	\item \textbf{失真度测试}：测试接收信号的失真程度
\end{enumerate}

\subsubsection{天线效应实验}
\begin{enumerate}[leftmargin=2cm]
	\item \textbf{天线长度影响}：测试不同长度天线对信号的影响
	\item \textbf{天线方向性}：测试天线的方向性特性
	\item \textbf{接地影响}：测试接地对天线性能的影响
	\item \textbf{天线匹配}：测试天线与信号源的阻抗匹配
\end{enumerate}

\subsection{实验注意事项}
进行无线实验时需要注意以下事项：

\begin{itemize}[leftmargin=2cm]
	\item \textbf{频率选择}：选择合法的实验频率，避免干扰正常广播
	\item \textbf{功率控制}：控制发射功率，避免过强信号干扰
	\item \textbf{天线安全}：确保天线安装安全，避免触电危险
	\item \textbf{环境选择}：选择开阔场地进行实验，减少干扰
	\item \textbf{法规遵守}：遵守无线电管理法规，仅在允许范围内实验
\end{itemize}

\subsection{实验数据分析}
实验过程中需要记录和分析以下数据：

\begin{itemize}[leftmargin=2cm]
	\item \textbf{信号强度}：记录不同距离的信号强度变化
	\item \textbf{频率稳定性}：记录信号频率的稳定性
	\item \textbf{调制质量}：评估调幅信号的质量
	\item \textbf{接收效果}：记录收音机的接收效果
\end{itemize}

通过以上实验，学生可以深入理解无线电波传播的基本原理，掌握无线通信的基本技术，为后续深入学习打下坚实基础。

\subsection{J24系列教学仪器概述}
J24系列是专门为中学物理教学设计的实验仪器系列，主要包括以下设备：

\subsubsection{J2464型教学信号源}
\begin{itemize}[leftmargin=2cm]
	\item \textbf{功能特点}：比J2465功能更丰富的教学信号源
	\item \textbf{输出波形}：提供多种波形输出，适合复杂实验
	\item \textbf{应用范围}：适用于高中物理电学实验教学
	\item \textbf{技术参数}：频率范围更宽，输出更稳定
\end{itemize}

\subsubsection{J2466型学生示波器}
\begin{itemize}[leftmargin=2cm]
	\item \textbf{功能特点}：专为学生设计的简易示波器
	\item \textbf{显示功能}：能够显示信号波形
	\item \textbf{教学应用}：用于观察电信号波形变化
	\item \textbf{操作特点}：操作简单，适合学生使用
\end{itemize}

\subsubsection{J2467型电学实验箱}
\begin{itemize}[leftmargin=2cm]
	\item \textbf{功能特点}：综合性电学实验装置
	\item \textbf{版本区别}：有木盒和塑料盒两种版本
	\item \textbf{实验内容}：包含多种基础电学实验
	\item \textbf{教学价值}：培养学生电路分析和设计能力
\end{itemize}

\subsubsection{J2468型其他教学仪器}
\begin{itemize}[leftmargin=2cm]
	\item \textbf{仪器类型}：可能包括其他专用教学设备
	\item \textbf{应用领域}：针对特定物理实验设计
	\item \textbf{教学目的}：满足不同层次的教学需求
\end{itemize}

\subsection{J24系列仪器的教学意义}
J24系列教学仪器具有重要的教育价值：

\begin{itemize}[leftmargin=2cm]
	\item \textbf{标准化设计}：统一的教学仪器标准
	\item \textbf{安全性保障}：符合教学安全要求
	\item \textbf{操作简便}：适合学生实验操作
	\item \textbf{教学配套}：与教材内容紧密结合
	\item \textbf{成本控制}：价格适中，适合学校采购
\end{itemize}

\subsection{J24系列仪器的发展历程}
J24系列教学仪器经历了以下发展阶段：

\begin{itemize}[leftmargin=2cm]
	\item \textbf{初期开发}：20世纪80年代开始研发
	\item \textbf{标准化进程}：逐步形成统一的技术标准
	\item \textbf{版本更新}：从木盒版本发展到塑料盒版本
	\item \textbf{技术改进}：不断优化性能和可靠性
	\item \textbf{教学应用}：在全国中学广泛使用
\end{itemize}

J24系列教学仪器作为中国特色的教学设备，在物理教育中发挥了重要作用，为培养学生科学素养和实践能力做出了积极贡献。

\chapter{微波测量仪器}
\section{微波功率计}
\section{微波频率计}
\section{微波网络分析仪}
\section{微波信号源}

\part{无线电测量系统}

\chapter{自动测试系统}
\section{自动测试系统的基本组成}
\section{GPIB总线系统}
\section{VXI总线系统}
\section{PXI总线系统}

\chapter{虚拟仪器系统}
\section{虚拟仪器的基本概念}
\section{LabVIEW虚拟仪器}
\section{基于PC的测量系统}
\section{软件定义无线电测试系统}

\chapter{远程测量系统}
\section{远程测量系统的组成}
\section{网络化测量系统}
\section{分布式测量系统}
\section{云计算测量系统}

\part{仪器校准与维护}

\chapter{仪器校准技术}
\section{校准的基本概念}
\section{校准方法和技术}
\section{校准标准与设备}
\section{校准周期与记录}

\chapter{仪器维护与管理}
\section{仪器的日常维护}
\section{故障诊断与维修}
\section{仪器管理信息系统}
\section{仪器寿命周期管理}

\chapter{测量不确定度评定}
\section{不确定度的基本概念}
\section{A类不确定度评定}
\section{B类不确定度评定}
\section{合成不确定度计算}

\part{无线电测量应用}

\chapter{通信系统测量}
\section{调制质量测量}
\section{误码率测量}
\section{频谱占用度测量}
\section{通信系统性能测试}

\chapter{雷达系统测量}
\section{雷达信号参数测量}
\section{雷达系统性能测试}
\section{雷达目标特性测量}
\section{雷达抗干扰性能测试}

\chapter{导航系统测量}
\section{卫星导航系统测试}
\section{惯性导航系统测试}
\section{无线电导航系统测试}
\section{组合导航系统测试}

\chapter{广播电视测量}
\section{广播发射机测试}
\section{电视信号质量测试}
\section{数字电视系统测试}
\section{广播电视覆盖测试}

\part{现代无线电测量技术}

\chapter{软件定义无线电测量}
\section{软件定义无线电的基本原理}
\section{SDR测量平台}
\section{软件定义测量技术}
\section{SDR在无线电测量中的应用}

\chapter{人工智能在无线电测量中的应用}
\section{机器学习在信号识别中的应用}
\section{深度学习在频谱分析中的应用}
\section{智能故障诊断技术}
\section{自适应测量技术}

\chapter{物联网测量技术}
\section{物联网设备测试}
\section{无线传感器网络测试}
\section{低功耗无线通信测试}
\section{物联网系统性能测试}

\chapter{5G及未来通信测量技术}
\section{5G信号测量技术}
\section{毫米波通信测试}
\section{大规模MIMO测试}
\section{未来通信系统测试展望}

\backmatter

\chapter{参考文献}

\chapter{附录}

\end{document}